\section{Weighted Dolbeault theory on the cyclic cover}\label{sec:weighted}

Lemma~\ref{compact:classes} supplies smooth end-framed line bundles
whose compact Chern classes satisfy $e^2=K_Y\cdot e<0$. We now give
each such bundle a holomorphic structure and a holomorphic section
approaching $1$ at the positive end and $0$ at the negative end.
The structure will approach the trivial one at any prescribed finite
exponential rate; that rate remains free throughout this section.

For a general compact datum we prove top-degree cohomology vanishing
and show that a holomorphic section is determined by its constant
coefficient at the positive end. Compact excision computes its weighted
Dolbeault index as $1+\tfrac12(e^2-K_Y\cdot e)$. For the classes above
this is one, so the vanishing and uniqueness statements give a unique
section after normalizing its positive coefficient to $1$. The
negative canonical degree of $e$ then forces its negative coefficient
to vanish.

The periodic input is Lemma~\ref{surface:characters}:
\begin{equation}\label{weighted:compact-vanishing}
 H^q(X,P_\lambda)=H^q(X,K_X\otimes P_\lambda)=0
 \quad(\lambda\ne1,\ 0\le q\le2),
 \qquad H^0(X,K_X\otimes P_\lambda)=0\quad(\lambda\in\CC^*).
\end{equation}
The height furnished by Lemma~\ref{surface:height} satisfies
$h\circ T=h+1$, $Lh=c>0$, and $|dh|\le C$. All its positive-order
derivatives are bounded. We fix the lifted Hermitian metric and its volume
form. The analysis of the ends uses these facts and the absence of
additive holomorphic modes established in
Lemma~\ref{surface:characters}.

\subsection{Spaces, adjoints, and Fredholm complexes}

For a compact line-bundle datum $E$ as defined in
Section~\ref{sec:compact-data}, choose a Hermitian metric and a smooth
connection agreeing with the standard ones in its prescribed end frames.
Changing these choices on a compact set gives equivalent Sobolev norms.
For $s\in\ZZ$, let $W^s$ denote the Sobolev space defined by a uniformly
locally finite lifted coordinate cover and partition of unity. For
negative $s$ use distributional duals of the positive-order local norms.
The lifted geometry has bounded derivatives, positive injectivity radius,
and uniformly finite overlap, so the usual localization, multiplication,
and elliptic estimates have constants independent of the lifted chart.

For real numbers $a,b$, choose a smooth function $\phi_{a,b}$ equal to
$ah$ on a sufficiently far positive tail and to $-bh=b|h|$ on a
sufficiently far negative tail. Its positive-order derivatives are
bounded. Set
\begin{equation}\label{weighted:norm}
 W^s_{a,b}(\Lambda^{p,q}\otimes E)
 =\{u:e^{-\phi_{a,b}}u\in W^s(\Lambda^{p,q}\otimes E)\},
 \qquad
 \|u\|_{s;a,b}=\|e^{-\phi_{a,b}}u\|_{W^s}.
\end{equation}
Any change in $\phi_{a,b}$ on a compact set gives the same space with an
equivalent norm. Positive $a$ or $b$ permits growth at the corresponding
end, and negative $a$ or $b$ requires decay. A single linear weight
$e^{-\rho h}$ has $(a,b)=(\rho,-\rho)$. We call a pair admissible when
$a\ne0$ and $b\ne0$.

Suppose for the moment that $E$ has an integrable operator $d_E$ of type
$(0,1)$ which, in its end trivializations, is
\begin{equation}\label{weighted:end-error}
 d_E=\dbar+A_E\wedge,
 \qquad |\nabla^j A_E|\le C_j e^{-\gamma |h|}
 \quad(j\ge0)
\end{equation}
for some $\gamma>0$. The same discussion applies to $E^{-1}$, to the
trivial bundle, and after tensoring by $K_Y$. For a fixed integer $s$,
the stepped Sobolev complex is the bounded complex
\begin{equation}\label{weighted:stepped}
 \mathcal C^{p,q}_{a,b,s}(E)=
 W^{s-q}_{a,b}(\Lambda^{p,q}\otimes E),
 \qquad d_E:\mathcal C^{p,q}_{a,b,s}(E)
       \longrightarrow\mathcal C^{p,q+1}_{a,b,s}(E),
 \quad 0\le q\le2.
\end{equation}
The decrease of Sobolev order with $q$ makes the differential bounded;
it will also allow the same theory to handle residue currents.
The cohomology will be denoted by $H^{p,q}_{a,b}(Y,E)$; its independence
of $s$ follows below. An alternative realization uses the maximal closed
operators $d_E$ in the equal-order weighted $L^2$ spaces. We shall verify
that this Hilbert complex has the same cohomology.

The analytic background used here is the compact elliptic Hodge theorem,
the local elliptic estimate (including its distributional Sobolev
versions), Rellich's compact embedding on relatively compact sets, and
Plancherel's theorem. On a compact manifold, the Hodge theorem for an
elliptic complex provides a finite-dimensional smooth harmonic space and
a Green operator of order $-2$, with
$I=P+d d^*G+d^*Gd$. The local elliptic estimate raises regularity by the
order of an elliptic operator when its right hand side has the
corresponding regularity. These results apply to smooth bundles and
smooth elliptic operators, with no K\"ahler hypothesis; for compact
Hermitian Dolbeault theory see
\cite[Chapter~VI, Theorems~(7.1)--(7.2)]{Demailly2012}. We give the
periodic and two-ended arguments, rather than import a theorem about a
different complex. The periodic-end Fredholm method goes back to Taubes
\cite[Lemma~4.3]{Taubes1987}; the de Rham formulation in
\cite[Proposition~2.1]{MRS2014} is useful background. The two-ended
Dolbeault argument needed here is proved below.

\begin{proposition}\label{weighted:fredholm}
For $p=0$ or $p=2$, every admissible pair $(a,b)$ and every $s\in\ZZ$,
the complex~\eqref{weighted:stepped} has closed ranges and
finite-dimensional cohomology. Its cohomology agrees with that of the
maximal weighted $L^2$ complex and is represented by forms which belong
to $W^k_{a,b}$ for every $k$. There are bounded homotopies of degree $-1$
which gain one derivative, modulo projections onto these harmonic
spaces. For the periodic trivial datum the single-linear-weight complex
with $\rho\ne0$ is acyclic, in every Sobolev realization.
The operator bounds may depend on the fixed rates and the datum.
\end{proposition}

\begin{proof}
First take a periodic complex on the whole cover. Conjugation by
$e^{-\rho h}$ changes its differential to
\[
 d_\rho=d+\rho\,\dbar h\wedge.
\]
Fourier transformation of the deck variable decomposes this complex as a
direct integral of the compact complexes on $X$ with multipliers
$e^{\rho+i\theta}$, $0\le\theta\le2\pi$. One can see both the
identification and the norm assertion on a compact fundamental region:
the sequence of its translates is an $\ell^2$ sequence with values in
local Sobolev spaces, and multiplication by $e^{-\rho h}$ changes the
norm only by bounded factors on that region. Plancherel's theorem gives
the direct integral. The endpoint gauges at $\theta=0,2\pi$ agree by
the unitary gauge $e^{2\pi i h}$; thus a global logarithm of the
multiplier is unnecessary.

By~\eqref{weighted:compact-vanishing}, every compact complex on this
circle is acyclic. Its Hodge operator $d_\rho+d_\rho^*$ is consequently
invertible. The compact elliptic estimate, the bounded inverse theorem,
and a finite covering of the parameter circle give a uniform inverse
from $L^2$ to $W^1$. For completeness, locally in the parameter an
already invertible operator remains invertible by a Neumann series;
the difference of two nearby operators is a small bounded map
$W^1\to L^2$. Applying the compact elliptic estimate to the inverse
gives the uniform higher-order bounds. Fourier inversion therefore gives
an inverse for the conjugated periodic Hodge operator, and the compact
homotopy identities integrate to an acyclic complex with a homotopy
gaining one derivative. Duality gives the same bounds at negative
Sobolev orders.

The adjoint used for a general weight is specifically the adjoint in
that weighted norm:
\begin{equation}\label{weighted:adjoint}
 d_{E,\phi}^*=e^{2\phi}d_E^*e^{-2\phi}.
\end{equation}
Thus $d_E+d_{E,\phi}^*$, conjugated by $e^{-\phi}$, is
$\widetilde d_E+\widetilde d_E^*$, where
$\widetilde d_E=e^{-\phi}d_Ee^\phi$. At the positive end its periodic
limit is the operator just considered with $\rho=a$; at the negative
end it is the one with $\rho=-b$. The difference has order zero and all
its coefficients tend to zero, by~\eqref{weighted:end-error}.

Here is the parametrix argument explicitly. On each tail use the inverse
of the corresponding whole-cover periodic limit, with nested cutoffs
equal to one sufficiently far along that tail. On a relatively compact
middle region use an interior elliptic parametrix. Compose their sum
with the conjugated Hodge operator on either side. Commutators with the
cutoffs are order-zero operators supported in compact collars. Following
or preceding the order-$-1$ parametrices they are compact, by Rellich's
theorem. Multiplication by a coefficient tending to zero is compact
from $W^1$ to $L^2$: truncate it to a compact set, apply Rellich there,
and bound the remaining operator norm by its supremum on the tails.
This treats the end errors as well. The resulting left and right
parametrices are inverses modulo compact operators. In particular the
Hodge operator is Fredholm from $W^1$ to $L^2$.

This construction also proves the useful estimate
\begin{equation}\label{weighted:fredholm-estimate}
 \|u\|_{1;a,b}\le C\bigl(
 \|(d_E+d_{E,\phi}^*)u\|_{0;a,b}
       +\|u\|_{L^2(C)}\bigr)
\end{equation}
for a fixed compact set $C$, and gives the analogous estimate for the
adjoint. The global first-order elliptic estimate shows that the
maximal domain of the conjugated Hodge operator is $W^1$. To justify
the assertion for an initially distributional element of that domain,
apply the local estimate on the uniform charts and sum; cutoffs in
$h/R$ then exhaust $Y$, with bounded commutators tending to zero.
The same cutoffs show that compactly supported smooth sections are a
core. The operator is therefore self-adjoint on $W^1$.

Because $d_E^2=0$, its square is the degree-preserving Hodge Laplacian.
Fredholmness and self-adjointness give a spectral gap away from its
finite-dimensional kernel. Let $P$ be the orthogonal projection onto
that kernel and let $G$ be the inverse of the Laplacian on its
orthogonal complement, zero on the kernel. Uniform elliptic estimates
give
\[
 G:W^t_{a,b}\longrightarrow W^{t+2}_{a,b}
 \quad(t\in\ZZ),\qquad
 Q=d_{E,\phi}^*G:W^t_{a,b}\longrightarrow W^{t+1}_{a,b}.
\]
For negative $t$, these follow also by transposing the positive-order
estimates after conjugation. Elements of $\operatorname{im}P$ have all
weighted Sobolev orders, so $P$ extends to all these spaces. Commutation
with the differential, first on smooth vectors and then by density,
gives
\begin{equation}\label{weighted:homotopy}
 I=P+d_EQ+Qd_E.
\end{equation}
On the stepped complex $Q$ maps degree $q$ into degree $q-1$ with
exactly the required gain of one derivative. A closed element has zero
class precisely when its projection under $P$ is zero. This proves
closed ranges, finite-dimensionality, and independence of the Sobolev
realization. It proves the same statements for the maximal $L^2$
complex, completing the proof.
\end{proof}

\begin{lemma}[Weighted duality]\label{weighted:duality}
For $p\in\{0,2\}$ and $0\leq q\leq2$, the integration pairing gives
a perfect pairing
\begin{equation}\label{weighted:serre}
 H^{p,q}_{a,b}(Y,E)\ \times\
 H^{2-p,2-q}_{-a,-b}(Y,E^{-1})\longrightarrow\CC.
\end{equation}
Here both pairs of rates are admissible. The statement concerns the
cohomology classes in any of the stepped realizations above; the pairing
is evaluated on their smooth weighted harmonic representatives.
\end{lemma}

\begin{proof}
Choose $\phi_{-a,-b}=-\phi_{a,b}$, so their exponential weight factors
are reciprocal. On weighted $L^2$
representatives the integrand is absolutely
integrable by Cauchy--Schwarz, since the two weights cancel. Integration
by parts is valid by inserting the cutoffs in $h/R$: their derivatives
are bounded and the $L^2$ tails tend to zero. Thus exact classes pair
trivially with closed classes. The Hermitian Hodge star and the bundle
metric give an antilinear map $J$ from $(p,q)$ forms with values in $E$
to complementary forms with values in $E^{-1}$, characterized by the
pointwise Hermitian inner product under wedge integration. The map
$u\mapsto e^{-2\phi}Ju$ takes the weighted harmonic equations to the
opposite-weight harmonic equations; this follows immediately by
transposing $d_E$ and using~\eqref{weighted:adjoint}. It is an
isomorphism, and the pairing of $u$ with its image is its weighted
squared norm, up to the harmless fixed degree sign. Harmonic
representatives therefore give a nondegenerate pairing. Both spaces
are finite-dimensional by Proposition~\ref{weighted:fredholm}.
At the cochain level, the dual of the order $s-q$ term has order $q-s$.
Writing the complementary degree as $2-q$ shows that the dual stepped
complex has base order $2-s$, since
$(2-s)-(2-q)=q-s$. Thus opposite weights also have this complementary
Sobolev shift; independence of the shift, already proved above, is what
allows the notation in~\eqref{weighted:serre} to suppress it.
\end{proof}

\subsection{Continuation in the multiplier and positive-tail uniqueness}

Fredholmness gives finite-dimensional kernels, but the index argument
requires more: scalar-type sections must have a single constant
coefficient at the positive end, and canonical sections must vanish
there. We obtain these conclusions by improving exponential rates.
First, the sign $Lh>0$ converts decay at every rate into vanishing on
the positive tail. We will then prove the required decay by continuing
the Fourier--Laurent transform in its multiplier.

\begin{lemma}[Positive-tail nondecay]\label{weighted:nondecay}
Let $V$ be a holomorphic line bundle on $\{h>H_0\}$, with a smooth
Hermitian metric of bounded Chern curvature. If a holomorphic section
$s$ satisfies, for every $N>0$,
\[
 |s(x)|\le C_Ne^{-Nh(x)}\quad(h(x)\text{ sufficiently large}),
\]
then $s$ vanishes on a positive tail. On every connected region on which
it is holomorphic and which meets that tail it consequently vanishes
identically.
\end{lemma}

\begin{proof}
The logarithm of the norm of a holomorphic section satisfies
$L\log|s|\ge-M$ in distributions, including at its zeros, for a fixed
$M$ depending on the curvature bound. In a holomorphic frame this is
the fact that the logarithm of the modulus of a holomorphic function is
plurisubharmonic, plus the bounded curvature contribution of the metric.
Consequently
$u=\log|s|+(M/c)h$ is subharmonic for $L$.
The chain rule and the bound on $dh$ give, for a sufficiently small
$\eta>0$,
\[
 L(e^{-\eta h})
 =e^{-\eta h}\bigl(-\eta c+\eta^2\,\sigma_L(dh,dh)\bigr)\le0,
\]
where $\sigma_L(dh,dh)$ is bounded and nonnegative.

Choose a regular value $h_0>H_0$ and a finite upper bound $C$ for $u$
on its compact level. Given $A>\max(1,-C)$, choose a regular $H>h_0$
so large that $u\le-A$ on $\{h=H\}$. Such an $H$ exists by the assumed
decay. The superharmonic function
\[
 v_A=-A+(A+C)e^{-\eta(h-h_0)}
\]
has value $C$ on the lower boundary and is at least $-A$ on the upper
boundary. The maximum principle on the compact slab gives $u\le v_A$.
It applies to the upper-semicontinuous subharmonic logarithm with its
value $-\infty$ at zeros, either directly or by local regularization.
For a fixed point of height greater than $h_0$, take $H$ beyond that
point and let $A\to\infty$. Since
$1-e^{-\eta(h-h_0)}>0$, we obtain $u=-\infty$ there. No connectedness
assumption on the level sets was used.
\end{proof}

Here and below, a section has \emph{finite exponential growth on a tail}
if it belongs there to some weighted Sobolev space of finite order and
finite real rate. Distributional holomorphic sections are smooth by
local elliptic regularity. The uniform local estimates then promote
this membership to every Sobolev order on a shorter tail, and to
pointwise bounds of the corresponding exponential size.

\begin{lemma}[Degree-zero continuation]\label{weighted:continuation}
Let $V_0$ be either $\OO_Y$ or $K_Y$ with its periodic holomorphic
operator $d$. Suppose on one tail that
\[
 (d+B\wedge)s=0,\qquad
 |\nabla^jB|\le C_j e^{-\gamma|h|},\quad j\ge0,
\]
where the perturbed operator is integrable and $\gamma>0$. If $s$ has
finite exponential growth, its allowed rate can be improved by any
amount strictly less than $\gamma$, except that a contour passing the
unit multiplier can contribute a constant when $V_0=\OO_Y$. There is
no degree-zero contribution for $V_0=K_Y$.

In the scalar case, after finitely many such improvements there is a
constant $c_\pm$ on the relevant end and an exponentially decaying
remainder, with decay in all derivatives. If that constant is zero,
the section has every finite exponential decay rate on that end. In
the canonical case it always has every finite exponential decay rate.
\end{lemma}

\begin{proof}
We transform a cut-off section to the compact base, continue the
resulting solutions through the exceptional multiplier, and recover
the improved rates by contour inversion.

\emph{The transformed equation.}
Consider first the positive tail and
choose a cutoff $\chi$ equal to one far along it and zero off a slightly
longer tail. In the periodic end frame, extend $u=\chi s$ by zero to
$Y$. Then
\begin{equation}\label{weighted:cutoff-equation}
 du=f=(d\chi)s-\chi Bs.
\end{equation}
If $s$ belongs to rate $r$, multiplication by $B$ puts the second term
at rate $r-\gamma$; the first term is compactly supported. Estimates
with any strictly intermediate rate avoid boundary-convergence issues.

For a positive-tail-supported form $v$, its Fourier--Laurent transform
is, initially in a right half-plane,
\begin{equation}\label{weighted:transform}
 \widehat v(z)(x)=
 \sum_{k\in\ZZ}e^{-z(h(x)+k)}(T^{k})^*v(x).
\end{equation}
The expression is invariant under $T$, hence a form on $X$. If $v$ has
rate $r$, Cauchy--Schwarz on the translates gives normal convergence
in compact-manifold Sobolev norms when $\operatorname{Re}z>r$.
The differentiated series converges on every smaller half-plane,
including any polynomial factors in $k$, so it is holomorphic there.
The transform of~\eqref{weighted:cutoff-equation} is
\begin{equation}\label{weighted:transformed-equation}
 d(z)\widehat u(z)=\widehat f(z),\qquad
 d(z)=d+z\,\dbar h\wedge.
\end{equation}
Its multiplier is $\lambda=e^z$; this also fixes our sign convention.
The source is holomorphic on $\operatorname{Re}z>r-\gamma$.
Moreover $d(z)\widehat f(z)=0$ there: the identity holds in the
initial half-plane because $d^2u=0$, and extends by the identity theorem.

\emph{Continuation on the compact base.}
For $e^z\ne1$, compact acyclicity gives a unique degree-zero solution
of~\eqref{weighted:transformed-equation}. It depends holomorphically on
$z$. To verify the latter assertion without choosing a nonholomorphic
family of Hodge inverses, fix $z_0$ where degree-zero cohomology vanishes.
Compact ellipticity makes
\[
 B_0=(d(z_0)^*d(z_0))^{-1}d(z_0)^*
\]
a bounded left inverse, from $W^t$ to $W^{t+1}$. Writing
$d(z)=d(z_0)+(z-z_0)A$, any existing solution obeys
\begin{equation}\label{weighted:left-inverse}
 \widehat u(z)=
 [I+(z-z_0)B_0A]^{-1}B_0\widehat f(z).
\end{equation}
The right side is holomorphic near $z_0$. In the canonical case
degree-zero cohomology also vanishes at $e^{z_0}=1$, so the same formula
continues the solution across that point. The full equation, not only
its image under $B_0$, holds by analytic continuation. Higher canonical
cohomology at this parameter creates no pole in this degree-zero
calculation.

In the scalar case the constants prevent such a left inverse at the
exceptional point. A local gauge reduces that point to $z_0=0$.
We first solve for the component orthogonal to the constants, then
determine the constant coefficient. Let $P_0$ be
projection onto the constants and use the compact left inverse $B_0$
with $B_0d=I-P_0$. Write a prospective solution as $c+v$, with $v$
orthogonal to the constants, and put
\[
 R(z)=(I+zB_0A)^{-1},\qquad
 v=R(z)B_0\widehat f-zcR(z)B_0A1.
\]
The inverses here act on the complement of the constants. Set
$w(z)=1-zR(z)B_0A1$; the candidate solution is
$R(z)B_0\widehat f+c w(z)$.
The class of $A1=\dbar h$ in $H^1(X,\OO_X)$ is nonzero. Indeed, if
$\dbar h=\dbar v$ for a function on $X$, then $h-v$ would be a
holomorphic function on $Y$ changing by $1$ under $T$, contrary to
Lemma~\ref{surface:characters}. Choose a continuous scalar
functional $\ell$ on degree-one forms which annihilates $\operatorname{im}d$
and for which $\ell(A1)\ne0$; pairing with the harmonic representative
of this nonzero class gives such a functional. Applying $\ell$ to
the transformed equation determines the candidate coefficient
\begin{equation}\label{weighted:scalar-pole}
 c(z)=
 \frac{\ell\bigl(\widehat f-d(z)R(z)B_0\widehat f\bigr)}
      {\ell(d(z)w(z))},
 \qquad
 \ell(d(z)w(z))=z\ell(A1)+O(z^2).
\end{equation}
This coefficient is meromorphic with at most a simple pole.
For every sufficiently small nonzero $z$, compact acyclicity already
supplies a unique solution of the full transformed equation. That
solution must satisfy both the complement equation and the displayed
scalar equation, so it equals our candidate. Thus the candidate solves
the full equation on the punctured neighborhood and continues it
meromorphically across zero. The residue
is a constant, since the pole part in the differential equation is
annihilated by $d$. This calculation supplies the required finite-order
meromorphic continuation without applying a Fredholm theorem to the
overdetermined map $d:W^{t+1}(\Lambda^{0,0})\to W^t(\Lambda^{0,1})$.

\emph{Recovering the end expansion.}
For a desired positive-end rate $r_1$ with $r-\gamma<r_1<r$, choose
a contour real part $\sigma\ne0$ such that
\[
 r-\gamma<\sigma<r_1<r.
\]
Start from any vertical line with real part greater than $r$.
Both lines lie in the source's domain of holomorphy, and the continued
solution has bounded Sobolev norm on each multiplier circle.
Moving the inversion contour to real part $\sigma$ gives, on a fixed
fundamental region,
\begin{equation}\label{weighted:inverse-transform}
 (T^k)^*u(x)=\frac1{2\pi i}
 \int_{\sigma-i\pi}^{\sigma+i\pi}
     e^{z(h(x)+k)}\widehat u(z)(x)\,dz,
\end{equation}
plus the residues crossed when moving from the initial line to $\sigma$.
Horizontal boundary terms cancel in the contour rectangle: the
transform changes by the gauge $e^{-2\pi i h}$ under $z\mapsto z+2\pi i$,
and the factor in~\eqref{weighted:inverse-transform} changes by its
inverse. This justifies the contour deformation. Cauchy's estimate on
the new line bounds its integral contribution on the $k$th translate
by $Ce^{\sigma k}$ in every fixed Sobolev norm. After applying the
target weight $e^{-r_1 k}$, its squared norms on the positive translates
are summable because
$r_1-\sigma>0$. The constants in~\eqref{weighted:scalar-pole} are the only
residues; there are none in the canonical case.

At the negative end use a cutoff supported in $\{h<H\}$. Its transform
converges first in a left half-plane. If the allowed negative-end rate
is $b$, the initial domain is $\operatorname{Re}z<-b$ and the source
extends to $\operatorname{Re}z<-b+\gamma$. Given $b_1$ with
$b-\gamma<b_1<b$, choose a new contour with
$-b_1<\sigma<-b+\gamma$ and $\sigma\ne0$. On the $(-l)$th translate
its integral is $O(e^{-\sigma l})$, whose squared norms become
summable after multiplication by $e^{-b_1l}$ because
$\sigma+b_1>0$. Thus moving the contour to the right improves the
negative-end rate; the residue discussion is unchanged.

It remains to check the iteration, since $\gamma$ is fixed. For a
homogeneous perturbed equation, each improvement in $s$ improves the
source $-Bs$ by the same amount. Choose, for example, improvements
smaller than $\gamma/3$ at each step, perturbing contour positions if
necessary to avoid real part zero. Canonical sections have no residue,
so iteration reaches every finite decay rate. Scalar sections can
cross zero once, giving $s=c_\pm+r$ with $r$ at a strictly negative
allowed rate. If $c_\pm=0$, the same homogeneous iteration continues
indefinitely. If $c_\pm\ne0$, then
\[
 (d+B\wedge)r=-B c_\pm;
\]
only a finite decay rate for this remainder is claimed. Uniform local
elliptic estimates give the stated decay of all its derivatives.
\end{proof}

\begin{remark}
The continuation calculation also identifies the possible homogeneous
modes. More generally, a pole of finite order at the unit multiplier
would give a polynomial-times-equivariant homogeneous mode.
Multiplication of the transform by a power of $e^z-1$ removes that
pole, so the corresponding section is killed by a power of $T^*-1$.
Such a nonconstant scalar mode would have a last nonconstant member
$v$ with $(T^*-1)v$ a nonzero constant, excluded by
Lemma~\ref{surface:characters}. A canonical generalized mode would
have a last nonzero invariant member, contradicting $H^0(X,K_X)=0$.
Thus the explicit simple-pole and removable-singularity calculations
agree with the absence of these generalized modes.
\end{remark}

\begin{proposition}[Tail asymptotics]\label{weighted:tail}
Let $E$ satisfy~\eqref{weighted:end-error}. A finite-growth holomorphic
section of $K_Y\otimes E$ or $K_Y\otimes E^{-1}$ on a positive tail
vanishes on a shorter positive tail. A finite-growth holomorphic section
of $E$ has, in the prescribed trivialization at either end on which it
is defined, the expansion
\[
 s=c_\pm+O(e^{-\delta |h|})
\]
for some $\delta>0$, with the same conclusion for every derivative of
the error. On the positive end $c_+=0$ implies $s=0$. Consequently the
positive-end coefficient is injective on global weighted holomorphic
sections of $E$.
\end{proposition}

\begin{proof}
Apply Lemma~\ref{weighted:continuation}. The Chern curvatures of the
metrics used here are bounded: those of the periodic bundles are
bounded, and the connection coefficients and their first derivatives in
the end frames are bounded by~\eqref{weighted:end-error}. Every
canonical section, and a scalar-type section with zero positive
coefficient, decays faster than every exponential on the positive
tail. Lemma~\ref{weighted:nondecay} makes it zero there. Unique
continuation for a holomorphic section on the connected $Y$ then gives
the global injectivity assertion. The negative-end expansion uses no
maximum principle at that end and makes no vanishing assertion there.
\end{proof}

Positive-tail uniqueness now controls the kernel of every compact datum.
To obtain such a datum holomorphically, we next solve the scalar
top-degree equation with arbitrarily rapid decay.

\subsection{Rapid scalar solutions and holomorphic compact data}

\begin{lemma}\label{weighted:scalar-top}
For the scalar periodic Dolbeault complex and every admissible pair,
\[
 H^{0,2}_{a,b}(Y,\OO)=0.
\]
In particular, if $f$ is a smooth compactly supported $(0,2)$-form and
$R>0$, there is a smooth $(0,1)$-form $\alpha$ with
$\dbar\alpha=f$ and
$\alpha\in W^k_{-R,-R}$ for every $k\in\ZZ$.
\end{lemma}

\begin{proof}
By Lemma~\ref{weighted:duality} the dual of the group in question is
$H^{2,0}_{-a,-b}(Y,\OO)$. Its elements are canonical sections with
finite exponential growth on the positive end; they vanish by
Proposition~\ref{weighted:tail}, applied to the periodic trivial datum.
Closed range and finite-dimensionality turn the dual vanishing into
the asserted vanishing. For $f$ use the homotopy $Q$ from
\eqref{weighted:homotopy} with weights $(-R,-R)$. The top-degree
harmonic projection is zero, so $\alpha=Qf$ solves the equation.
Elliptic regularity and the higher-order bounds for $Q$ give all the
asserted weighted Sobolev orders.
\end{proof}

\begin{proposition}\label{weighted:holomorphic}
Every compact line-bundle datum admits an integrable $(0,1)$-connection
satisfying~\eqref{weighted:end-error} with any prescribed finite
$\gamma>0$. For every such structure and every admissible pair,
\begin{equation}\label{weighted:E-top}
 H^{0,2}_{a,b}(Y,E)=0.
\end{equation}
For every $a,b>0$, global weighted holomorphic
sections have the two constant end coefficients and the exponentially
decaying errors described in Proposition~\ref{weighted:tail}; the map
$s\mapsto c_+$ is injective.
\end{proposition}

\begin{proof}
Choose a smooth $(0,1)$-connection $d_0$ on $E$ which is standard in the
given end trivializations. A partition of unity constructs one, because
the space of such connections is affine. Its curvature $d_0^2$ is a
compactly supported scalar $(0,2)$-form $\kappa$: the endomorphism
bundle of a line bundle is canonically trivial. In complex dimension
two it is automatically $\dbar$-closed. Apply
Lemma~\ref{weighted:scalar-top}, already proved solely for the trivial
periodic bundle, to solve
\[
 \dbar\alpha=-\kappa
\]
with weights $(-R,-R)$, where $R>\gamma$. Define
$d_E=d_0+\alpha\wedge$. In a local frame its square is
\[
 d_E^2=\kappa+\dbar\alpha+\alpha\wedge\alpha=0.
\]
The last term is zero because $\alpha$ is scalar-valued and has degree
one. The local integrability theorem for a $(0,1)$-connection therefore
makes $d_E$ a holomorphic line-bundle structure. In rank one this
theorem is just the local Dolbeault lemma applied to its closed
connection form, followed by multiplication of the frame by the
exponential of a local primitive. Weighted Sobolev embedding on the
uniform charts gives
$|\nabla^j\alpha|\le C_j e^{-R|h|}$ on the tails, hence the prescribed
bound with $\gamma$. The finitely many middle charts impose no end
condition.

Now weighted duality identifies the dual of
$H^{0,2}_{a,b}(Y,E)$ with
$H^{2,0}_{-a,-b}(Y,E^{-1})$. The latter vanishes by
Proposition~\ref{weighted:tail}. This proves~\eqref{weighted:E-top}.
The section assertions follow from the same proposition, whose
finite-growth hypothesis permits every positive finite rate. The scalar
top-degree vanishing used to construct $d_E$ preceded, and did not
assume, the existence of this holomorphic structure.
\end{proof}

The preceding solution makes the line-bundle data integrable with a
prescribed end rate. Their top-degree cohomology is zero, and the
positive-end coefficient bounds the dimension of their section spaces
by one. The index calculation will show that this bound is attained
for the classes in Lemma~\ref{compact:classes}.

\subsection{The index and compact excision}

\begin{lemma}\label{weighted:base-index}
For the trivial datum and $a,b>0$,
\[
 H^{0,0}_{a,b}(Y,\OO)=\CC,
 \qquad H^{0,1}_{a,b}(Y,\OO)=H^{0,2}_{a,b}(Y,\OO)=0.
\]
In particular its weighted Dolbeault index is one.
\end{lemma}

\begin{proof}
A global weighted holomorphic function has a positive-end constant
coefficient. Subtract that constant; the resulting holomorphic function
has zero positive coefficient and vanishes by
Proposition~\ref{weighted:tail}. Constants belong to the space because
both end rates are positive. The top-degree assertion is
Lemma~\ref{weighted:scalar-top}.

Here are the details in degree one. Represent a class by a smooth
weighted harmonic $(0,1)$-form $u$, so all its weighted Sobolev norms
are finite. Choose $\chi_++\chi_-=1$ with $\chi_+$ zero far down the
negative end and one far up the positive end. The compactly supported
$(0,2)$-form
$g=\dbar(\chi_+u)=-\dbar(\chi_-u)$ admits a primitive $v$ decaying
at both ends at any prescribed finite rate, by
Lemma~\ref{weighted:scalar-top}. Choose that rate larger than $a+b+1$.
Then
\[
 u_+=\chi_+u-v,\qquad u_-=\chi_-u+v
\]
are closed, with $u=u_++u_-$. The first belongs to the single-linear
space $(a,-a)$, because its negative tail is $-v$; the second belongs
to the single-linear space $(-b,b)$, because its positive tail is $v$.
Both complexes are acyclic by Proposition~\ref{weighted:fredholm}.
Choose $\dbar f_+=u_+$ and $\dbar f_-=u_-$ in those respective spaces.
Each of these spaces embeds into the two-growth-end space $(a,b)$:
one end has the same weight and the other has a stronger decay
condition. Thus $u=\dbar(f_++f_-)$ in the original complex.
\end{proof}

\begin{proposition}\label{weighted:index}
Let $E$ be a compact datum supported in a cut-open period of a finite
cyclic quotient $X_n=Y/T^n$, with its trivializations fixed in collars
of the cut. Equip $E$ with an integrable operator $d_E$ satisfying
\eqref{weighted:end-error}, and let $e$ be its compact support Chern
class. For $a,b>0$ its weighted Dolbeault index is
\begin{equation}\label{weighted:index-formula}
 \sum_{q=0}^2(-1)^q\dim H^{0,q}_{a,b}(Y,E)
 =1+\frac12\bigl(e^2-K_Y\cdot e\bigr).
\end{equation}
The products on the right are compact support intersection pairings.
In particular, if $e^2=K_Y\cdot e$, the section space is
one-dimensional, its positive-end coefficient is nonzero on every
nonzero section, and $H^{0,1}_{a,b}(Y,E)=0$.
\end{proposition}

\begin{proof}
The even-to-odd operator
\[
 D_E^+=d_E+d_{E,\phi}^*:
 W^1_{a,b}(\Lambda^{0,0}E\oplus\Lambda^{0,2}E)
 \longrightarrow L^2_{a,b}(\Lambda^{0,1}E)
\]
has index equal to the Euler characteristic in the statement, by the
Hodge decomposition proved above. Replace $d_E$ by any smooth
$(0,1)$-connection $d_0$ standard at the ends. The difference, after
conjugation, is multiplication by exponentially decaying coefficients,
and is compact from $W^1$ to $L^2$. The index is unchanged. We do not
require $d_0^2=0$ for this replacement: $D_0^+$ is still a first-order
elliptic operator with the same principal symbol and the same invertible
end limits.

Transfer the datum and this connection to $X_n$, making them trivial and
standard across the cut collars. Denote the resulting smooth bundle
by $E_n$ and its elliptic operator by $D_{E_n}^+$. Modify
$\phi_{a,b}$ on a compact set so that it is zero on a neighborhood of
the entire transferred piece and its collars; this does not change
the index. The following excision uses the same compact cut-and-paste
principle as \cite[Proposition~3.1]{MRS2014}; we give the operator
comparison in the present setting:
\begin{equation}\label{weighted:excision}
 \ind D_{E,Y}^++\ind D_{\OO,X_n}^+
 =\ind D_{\OO,Y}^++\ind D_{E_n,X_n}^+.
\end{equation}
To see it directly, compare these two direct sums. On the common
compact interior interchange the two copies, carrying the $E$ copy
to the transferred $E_n$ copy and the trivial copy to the trivial copy.
Outside the identified region use the identity. In each collar both
bundles are trivialized; interpolate between interchange and identity
using the real rotation matrix
\[
 \begin{pmatrix}\cos\theta&-\sin\theta\\
                 \sin\theta&\cos\theta\end{pmatrix},
 \qquad 0\le\theta\le\pi/2.
\]
Use the same rotation on the even and odd bundles. This defines bounded
invertible maps on the direct-sum domains and ranges. On the interior
the operators are interchanged exactly. On the collars their principal
symbols are the identical Dolbeault symbols, which commute with the
rotation; differentiating the rotation only adds order-zero terms.
Outside the compact comparison region they agree. The conjugated
operators therefore differ by compactly supported order-zero terms,
which are compact from $W^1$ to $L^2$. Their indices agree, proving
\eqref{weighted:excision}.

The transferred connection on $X_n$ need not be integrable. To compute
its index nevertheless by ordinary Hirzebruch--Riemann--Roch, recall
from Lemma~\ref{surface:covers} that $H^2(X_n,\OO_{X_n})=0$. The
exponential sequence then shows that every integral class in $H^2(X_n;\ZZ)$
is the first Chern class of a holomorphic line bundle. Smooth complex
line bundles are classified by their first Chern class, so $E_n$ admits
some holomorphic structure. Replacing the transferred connection by
its Dolbeault operator changes the closed elliptic operator only in
order zero and hence preserves its index. Compact-complex Riemann--Roch
\cite[\S4]{AtiyahSingerIII1968} gives
\[
 \ind D_{E_n,X_n}^+-\ind D_{\OO,X_n}^+
 =\chi(X_n,E_n)-\chi(X_n,\OO_{X_n})
 =\frac12\bigl(c_1(E_n)^2-K_{X_n}\cdot c_1(E_n)\bigr).
\]
These products equal $e^2$ and $K_Y\cdot e$, since the class was
transferred in its fixed trivialized collars. Combining
\eqref{weighted:excision} with Lemma~\ref{weighted:base-index} proves
\eqref{weighted:index-formula}.

When the correction term is zero, Proposition~\ref{weighted:holomorphic}
gives $H^{0,2}=0$ and Proposition~\ref{weighted:tail} gives
$\dim H^{0,0}\le1$. The index equality then reads
$\dim H^{0,0}-\dim H^{0,1}=1$. Both dimensions are nonnegative,
so they are respectively $1$ and $0$, with the asserted nonzero
positive coefficient.
\end{proof}

\subsection{Sections with prescribed end limits}

For the compact classes constructed earlier the index correction
vanishes, so a section with positive coefficient $1$ now exists.
We finish the construction by determining its negative coefficient
and identifying the class of its divisor.

\begin{proposition}[A section with different end limits]
\label{compact:section}
Let $E$ be one of the data in Lemma~\ref{compact:classes}, and write
$e^2=K_Y\cdot e=-m$ with $m>0$. Given any finite $\gamma>0$, $E$ admits
a holomorphic structure whose Dolbeault operator, in its fixed end
frames, differs from the trivial operator by $O(e^{-\gamma|h|})$ in
every derivative. For this structure there is a holomorphic section $s$
such that, in those frames,
\begin{equation}\label{compact:end-limits}
 s\longrightarrow1\quad(h\to+\infty),\qquad
 s\longrightarrow0\quad(h\to-\infty),
\end{equation}
with exponentially decaying errors. Its effective divisor $D=(s=0)$
is contained in $\{h\leq H\}$ for some $H$, is locally finite, and
represents the ordinary Chern class $j(e)$, or equivalently its dual
locally finite homology class.
\end{proposition}

\begin{proof}
Proposition~\ref{weighted:holomorphic} supplies the stated holomorphic
structure with the prescribed finite rate $\gamma$. Choose any
$\eps>0$ and use the pair $(\eps,\eps)$.
Proposition~\ref{weighted:index} gives
\[
 \ind\dbar_E
 =1+\frac12(e^2-K_Y\cdot e)=1.
\]
By the weighted degree-two vanishing and the positive-end uniqueness
in Propositions~\ref{weighted:holomorphic} and~\ref{weighted:tail},
$H^{0,2}_{\eps,\eps}(Y,E)=0$ and
the positive leading coefficient is an injective map
\[
 H^{0,0}_{\eps,\eps}(Y,E)\longrightarrow\CC.
\]
The index identity therefore implies that this degree-zero space is
one-dimensional (and the degree-one space is zero). Take its nonzero
section and normalize the positive leading coefficient to one.
Proposition~\ref{weighted:tail} gives a constant negative leading coefficient
$c_-$ and exponentially small errors at both ends.

Suppose $c_-\neq0$. Uniform convergence to nonzero constants would make
$s$ nowhere zero outside a compact height slab. Its zero divisor $D$
would then be compact. Moreover its compact divisor class would equal
the particular relative class $e$, not just its ordinary image. To see
the role of the framings, write $s=f_\pm\tau_\pm$ on the ends in the
fixed frames. Far enough out, $f_+$ lies in a disk about $1$ that avoids
zero and $f_-$ lies in a disk about $c_-$ that avoids zero. Straight-line
homotopies in these disks deform $f_\pm$ through nowhere-zero maps to
their limiting constants. Each nonzero constant is in turn homotopic
to $1$ in $\CC^*$. Thus the trivializations provided by $s$ on the
ends are homotopic to the prescribed ones. The relative zero-divisor
formula for $c_1$ consequently identifies the compact divisor class
with $e$.

The restriction of $\pi$ to the compact divisor is proper. Its pushforward
$\pi_*D$ is an effective divisor on $X$, with positive integral
multiplicities equal to the corresponding generic mapping degrees.
Since $K_Y=\pi^*K_X$, the projection formula gives
\[
 K_X\cdot\pi_*D=K_Y\cdot D=K_Y\cdot e=-m<0.
\]
This contradicts Lemma~\ref{surface:degree}; if the divisor were empty,
the left side would instead be zero and give the same contradiction.
Therefore $c_-=0$, proving \eqref{compact:end-limits}.

The positive limit excludes zeros for all sufficiently large $h$.
The zero set of a nonzero holomorphic section is a locally finite
effective divisor. The local divisor construction identifies
$E\simeq\OO_Y(D)$, so $[D]=c_1(E)=j(e)$ in ordinary cohomology;
the corresponding Chern-current identity is the local form of
\cite[Chapter~V, (13.2)]{Demailly2012}.
Under Poincar\'e duality for the oriented noncompact four-manifold $Y$,
this says that its locally finite fundamental cycle is the dual of
$j(e)$. In particular, if $\alpha$ is any compactly supported closed
two-form, then
\begin{equation}\label{compact:divisor-pairing}
 \int_Y\alpha\wedge j(e)=\int_D\alpha.
\end{equation}
The right side is well-defined by local finiteness. No assertion that
$D$ is a compact divisor is being made here.
\end{proof}

The rate $\gamma$ in Proposition~\ref{compact:section} is free. The
next sections first obtain an area exponent $A$ from the fixed lifted
coordinate geometry, and then use this proposition with
$\gamma>2A+10$. Thus the choice of that exponent imposes no additional
condition on the compact classes constructed above.
